\section{Evidence, Receipts, and Gates}
\label{sec:sources-wasm4pm-evidence}

\subsection{Tests}
\label{sec:sources-wasm4pm-evidence-tests}

All 31 tests pass as of 2026-06-01, with a total runtime of 10.60 seconds across four test
suites: unit tests (inline within source modules), integration tests
(\texttt{tests/integration\_tests.rs}), end-to-end tests (\texttt{tests/e2e\_tests.rs}), and
weaver integration tests (\texttt{tests/weaver\_integration\_tests.rs}).

The conformance module (\texttt{src/conformance.rs}) carries six inline unit tests:
\texttt{test\_conformance\_verdict\_fitness} (verifies \texttt{ConformanceVerdict::fitness\_score}
returns correct values for all three variants);
\texttt{test\_conformance\_verdicts\_aggregation} (verifies aggregate fitness computation across
multiple cases); \texttt{test\_alignment\_fitness} (verifies the van der Aalst fitness equation
$f = 0.5(1 - m/c) + 0.5(1 - r/p)$ with known token counts);
\texttt{test\_astar\_alignment\_correctness\_and\_performance} (verifies A* search in
synchronous product space reaches correct alignment within the 5000-iteration cap);
\texttt{test\_compile\_time\_metrics} (verifies \texttt{Between01<NUM,DEN>} const generic
compilation with valid and invalid bounds); \texttt{test\_runtime\_metric\_validation} (verifies
\texttt{RuntimeBetween01} rejects NaN, $\pm\infty$, and out-of-bounds \texttt{f64} values).

The mining module (\texttt{src/mining/mod.rs}) carries seven inline unit tests:
\texttt{test\_alpha\_witness\_lattice\_bottom} (verifies \texttt{AlphaWitness::bottom()} returns
empty sets); \texttt{test\_inductive\_witness\_lattice\_top} (verifies
\texttt{InductiveWitness::top()} returns saturated counts); \texttt{test\_heuristics\_witness\_join}
(verifies lattice join is monotone); \texttt{test\_petri\_net\_serialization};
\texttt{test\_process\_tree\_serialization}; \texttt{test\_alpha\_miner\_discovery} (end-to-end
alpha miner from event log to \texttt{Evidence<ProcessModel::Net,~Admitted,~AlphaWitness>});
\texttt{test\_dfg\_miner\_discovery}. The graduation module (\texttt{src/graduation.rs}) carries
four inline tests: \texttt{test\_valid\_intake}, \texttt{test\_ungrounded\_intake},
\texttt{test\_reason\_mismatch}, and \texttt{test\_invalid\_witness\_state}, verifying all four
intake pathways of the \texttt{GraduateToWasm4pm} trait.

\subsection{Receipts}
\label{sec:sources-wasm4pm-evidence-receipts}

The primary receipt artifact is \texttt{MINING\_RENDER\_RECEIPT.md}: Receipt ID
\texttt{MINING\_AUTHORITY\_RENDER\_001}, Authority: Mining Authority wasm4pm v30.1.2, Timestamp:
2026-06-01T00:00:00Z, Status: SEALED. The sealed artifact is \texttt{src/mining/mod.rs} with
SHA-256 content hash
\texttt{08a067d1ee19ea67150c194e9a2db7d86dfd994223b922de7e8606f45fbdf8e5}. The receipt documents
six manufacturing steps: (1) extraction of mining authority specifications from
\texttt{mining-authority-map.md}; (2) template application from
\texttt{wasm4pm-compat/compat/templates/mining/}; (3) artifact materialization to
\texttt{src/mining/mod.rs}; (4) zero-warning compilation verification (\texttt{cargo build
-{}-release}, 1.73s); (5) \texttt{Evidence<T,~State,~W>} type correctness audit across all four
witness types; (6) authority sealing with receipt signature.

The \texttt{Evidence<T,~Admitted,~W>} carrier type carries a \texttt{Blake3Hash([u8;~32])} field
in every instance. As documented in the open questions (GAP from the provenancce chain analysis),
all current implementations initialize this field to \texttt{Blake3Hash([0u8;~32])}: the
cryptographic infrastructure is structurally present but hash computation over algorithm inputs is
not yet wired. This is a documented open gap, not a silent omission.

\subsection{Checkpoints}
\label{sec:sources-wasm4pm-evidence-checkpoints}

The \texttt{research-verdict.md} (dated 2026-05-31, Authority: Execution Agent, Classification:
Doctoral-Level Authority Specification) serves as the checkpoint document for this project. It
classifies four execution authorities to doctoral-level precision and issues a GO verdict for
engineering. The document records specification completeness for all four authorities (Mining,
Conformance, Replay, Lifecycle), documents all proof obligations, FFI boundaries, error handling,
and implementation guidance. The verdict confidence is stated as ``DOCTORAL THESIS (99\%
confidence in specification completeness).'' The \texttt{EXECUTION\_AUTHORITY\_ATLAS.md}
(dated 2026-05-31, Agent: E3, Status: RESEARCH PASS COMPLETE) serves as the audit checkpoint for
the live \texttt{\textasciitilde/wasm4pm/} codebase, confirming the ten gaps and the critical
GAP\_001 finding.

\subsection{ALIVE/PARTIAL/BLOCKED Status}
\label{sec:sources-wasm4pm-evidence-status}

\textbf{Status: ALIVE.} The status is ALIVE because all three required conditions are met: (1)
the receipt is present and sealed (\texttt{MINING\_RENDER\_RECEIPT.md}, status SEALED,
2026-06-01T00:00:00Z); (2) all 31 tests pass (unit + integration + end-to-end + weaver,
10.60s runtime); (3) the research verdict checkpoint (2026-05-31) issues a GO verdict with
doctoral-level specification completeness for all four execution authorities. The ten open gaps
(GAP\_001 through GAP\_010) documented in \texttt{GAP\_ANALYSIS.md} represent known open
research findings. GAP\_001 (zero \texttt{wasm4pm-compat} consumption in the live
\texttt{\textasciitilde/wasm4pm/} codebase) is the most critical; it does not block the ALIVE
verdict for \texttt{sources/wasm4pm} because the local \texttt{src/graduation.rs} bridges the
type-law seam on this side of the boundary. The ALIVE verdict reflects the state of this
manufactured codebase; the relationship to the live \texttt{\textasciitilde/wasm4pm/} codebase
is the primary open research question.
